% Nouveau EGA V (cote 25): archive cote 25, site page 123 (PDF page 124).
% typescript. Draft compared: none.
% AI-transcribed by Claude Opus 5.5 (claude-opus-5-5, adaptive thinking, effort high), Message Batches API, batch ega25b (msgbatch_01A7n5YuXubUPNK1zsepi8RG), 2026-11-13; not yet checked by eye.
\scanpage{25/123}
\mspage{26}

\var{\over{7}{4}}{4}{writing over}{ega25_25_123-v0001}. \uline{Systèmes linéaires de diviseurs et Modules inversibles}.

Utilisant le résultat du nº1, nous allons donner une description
complète des systèmes linéaires sur X en termes des faisceaux inversibles
sur X. \var{\inserted{Nous pouvons supposer évidemment $\check{P}\to S$ surjectif. Alors $X\times_S\check{P}\to \unsure{\check{P}}$ l'étant.}}{Nous pouvons supposer évidemment $\check{P}\to S$ surjectif. Alors $X\times_S\check{P}\to \unsure{\check{P}}$ l'étant.}{insertion}{ega25_25_123-v0002} \var{\deleted{En effet,}}{}{deletion}{ega25_25_123-v0003} d'après les généralités du \var{\over{nº2}{\unsure{\S 20.3}}}{\unsure{\S 20.3}}{writing over}{ega25_25_123-v0004} la donnée d'un
diviseur \var{\deleted{de Cartier}}{}{deletion}{ega25_25_123-v0005} D sur $X\times_S\check{P}$ \var{\xout{(fibré projectif sur X)}}{}{typewriter strike-out}{ega25_25_123-v0006} revient à la
donnée d'un Module inversible $\underline{N}$ sur $X\times_S\check{P}$ , et d'une section régulière
$\phi$ de ce dernier. \var{\xout{\unsure{En ver}}}{}{typewriter strike-out}{ega25_25_123-v0007} L'hypothèse que D est un\var{\xout{e}}{}{typewriter strike-out}{ega25_25_123-v0008} système linéaire de
diviseurs sur X paramétrée par $\check{P}$ s'exprime alors par les deux conditions
1º) les $\phi_t$ ($t\in\check{P}$) induits par $\phi$ sur les fibres de $Xs_S\check{P}$ au dessus de $\check{P}$
sont réguliers (ce qui entraine que $\phi$ est régulier), et 2º) les $\underline{N}_x$ ($x\in X$)
induits par $\underline{N}$ \var{\xout{xxxx}}{}{typewriter strike-out}{ega25_25_123-v0009} sur les fibres des $X\times_S\check{P}$ sur X sont \uline{de degré} 1 .
Or la donnée d'un $\underline{N}$ inversible sur \unsure{le} fibré projectif $X\times_S\check{P}$ sur X,
satisfaisant à la condition 2º) ci-dessus, équivaut en vertu du nº1
à la donnée d'un Module inversible $\underline{L}$ sur X, $\underline{N}$ étant déterminé en fonction
de $\underline{L}$ par
\[ \underline{N}\simeq\underline{L}\otimes_{\underline{O}_{\unsure{S|X}}}\underline{O}_{\check{P}}(1) \quad , \]
et $\underline{L}$ étant d'ailleurs déterminé en termes de $\underline{N}$ par $\underline{L}\simeq pr_{1*}(\underline{N}(-1))$,
où (-1) désigne \var{\inserted{\unsure{la|le}}}{\unsure{la|le}}{insertion}{ega25_25_123-v0010} tensorisation par $\underline{O}_{\check{P}}(-1)$ sur $\underline{O}_S$ . La donnée de $\phi$
revient à la donnée d'une section de $\underline{L}\otimes_{\underline{O}_S}\underline{O}_{\check{P}}(1)$ , i.e. d'une section
\var{\margin{$\underline{O}_{X\times_S\check{P}}(1)$}}{\placedmargin{$\underline{O}_{X\times_S\check{P}}(1)$}}{marginal note placed}{ega25_25_123-v0011}
de $pr_{1*}(\underline{L}\otimes_{\underline{O}_X}\underline{O}_{\hand{X\times_S\check{P}}}(1))$ = $\underline{L}\otimes_{\underline{O}_S}\unsure{pr}_*(\underline{O}_{\hand{\check{P}}}(1))$, \var{\deleted{\var{\xout{xx$\underline{L}$x$\otimes$xx$\underline{E}$xxxi.e.xxd'unxhomomorphisme}}{}{typewriter strike-out}{ega25_25_123-v0013}}}{}{deletion}{ega25_25_123-v0012}
\var{\deleted{\var{\xout{$\underline{E}$x$\underline{E}$xxxx$\otimes$}}{}{typewriter strike-out}{ega25_25_123-v0015}}}{}{deletion}{ega25_25_123-v0014} or en vertu de III 2 on a $\unsure{pr}_*(\underline{O}_{X\times_S\check{P}}(1))\simeq\check{\underline{E}}_X$ , donc
la donnée de $\phi$ équivaut à \var{\xout{xxxx}}{}{typewriter strike-out}{ega25_25_123-v0016} la donnée d'un homomorphisme $g^*(\underline{E})\to\underline{L}$,
où ce qui revient au même, d'un morphisme \var{\inserted{u:}}{u:}{insertion}{ega25_25_123-v0017} $\underline{E}\to g_*(\underline{L})$ (NB $g:X\to S$ est
la projection canonique). Il reste à exprimer la condition 1º) plus
haut en termes de u. Comme les constructions faites commutent avec le
changement de base, il suffit d'exprimer cette condition fibre par fibre\hand{,}
et compte tenu que les points de $\check{P}$ à valeurs dans une extension k de
k(s) correspondent exactement aux droites de $E(s)\otimes_{k(s)}k$, cette

\mspage{123}
